\section{The equality case}
\label{sec:equality}

We show that equality in the grading bound forces the projective-space
product.  The nonzero product from
Proposition~\ref{prop:independent-correspondences} first gives a chain
of actual minimum-degree rational curves.  A dimension estimate then
turns the corresponding families into covering families.

We use $\operatorname{RatCurves}^{\mathrm n}(X)$ for the normalized
parameter space of rational curves, obtained from maps
$\mathbb P^1\to X$ birational onto their images by reparametrization.
In this section a family $V$ is an irreducible component of that space.
It is \emph{unsplit} if it is proper, and \emph{covering} if its
universal evaluation map dominates $X$.  Its anticanonical degree is
$-K_X\cdot V$, the common degree of its members.  For unsplit families
$V^1,\ldots,V^k$ and a point $x\in X$, write
$\operatorname{Locus}(V^1,\ldots,V^k)_x$ for the set of endpoints $y$
of chains $C_1,\ldots,C_k$ with $C_j\in V^j$, $x\in C_1$,
$C_j\cap C_{j+1}\ne\varnothing$, and $y\in C_k$.

We will use the following two established results, with their family
hypotheses stated explicitly.

\begin{theorem}[The ordered-chain bound]
\label{thm:ordered-chain-bound}
Let $Y$ be a smooth complex projective variety.  Let
$V^1,\ldots,V^k$ be proper irreducible components of
$\operatorname{RatCurves}^{\mathrm n}(Y)$ whose numerical classes are
linearly independent.  For every $y\in Y$, the ordered endpoint locus
is either empty or satisfies
\[
 \dim\operatorname{Locus}(V^1,\ldots,V^k)_y
 \ge\sum_{j=1}^k(-K_Y\cdot V^j-1).
\]
\end{theorem}

This is the theorem of Bonavero--Casagrande--Debarre--Druel
\cite[Theorem~5.2]{BCDD2003}, with the irreducible-component hypothesis
specified in their erratum~\cite{BCDDErratum}.  Taking arbitrary
proper subfamilies would not suffice.

\begin{theorem}[Occhetta's product characterization]
\label{thm:occhetta-product}
Let $Y$ be a smooth complex projective variety of dimension $n$.
Suppose $Y$ has numerically independent unsplit covering families
$V^1,\ldots,V^k$ with anticanonical degrees $n_1+1,\ldots,n_k+1$,
where each $n_j$ is positive and $\sum_jn_j=n$.  Then
$Y\cong\mathbb P^{n_1}\times\cdots\times\mathbb P^{n_k}$.
\end{theorem}

We use Theorem~\ref{thm:occhetta-product} in the form of
\cite[Theorem~1.1]{Occhetta2006}.  The next lemma verifies that the curve
families needed for these results are supplied by our correspondences.

\begin{lemma}[Minimum-degree incidence]
\label{lem:minimum-degree-incidence}
Let $\gamma_1,\ldots,\gamma_k$ be linearly independent numerical
stable-map degrees on $X$, each of anticanonical degree $\iota$.
If $S_{\gamma_1}\cdots S_{\gamma_k}\ne0$, there exist numerically
independent unsplit families $V^1,\ldots,V^k$, each of degree $\iota$,
and a nonempty chain in some ordering of those families.
\end{lemma}

\begin{proof}
For a numerical degree $\gamma$, let
$\overline{\mathcal M}_{0,2}(X,\gamma)$ denote the finite disjoint
union of stable-map stacks over actual curve classes of that degree.
Let $Z_\gamma\subset X\times X$ be the union of the images of their
two-evaluation maps.
It is closed by properness.  The correspondence inducing $S_\gamma$
is the pushforward of its virtual class and is supported on
$Z_\gamma$.  The composition of these correspondences is computed
on $X^{k+1}$ by imposing the consecutive-pair incidence conditions
and projecting to the two endpoints.  Refined intersection with the
diagonals is supported on the incidence fiber product.  If that fiber
product were empty, the resulting correspondence, and hence its
action on cohomology, would be zero.  Thus the nonzero product gives
a sequence of two-pointed stable maps with matching evaluations.
Choose their order according to composition, reversing the written
list of degrees if necessary.  This support argument does not require
effectivity of the virtual classes.

Every nonconstant component of a genus-zero stable map has an
irreducible rational image of anticanonical degree at least $\iota$.
Its contribution to the total degree is that image degree multiplied
by its covering multiplicity.  A stable map of total degree $\iota$
therefore has exactly one nonconstant component, and that component
is birational onto its image.  The images in our sequence are thus
irreducible rational curves $C_1,\ldots,C_k$ in the specified
numerical classes, with consecutive intersections nonempty.
Contracted components carrying markings do not affect this conclusion.

Choose an irreducible component $V^j$ of
$\operatorname{RatCurves}^{\mathrm n}(X)$ containing $C_j$.
Its members retain degree $\iota$.  To check properness, take stable
limits after forgetting markings.  Such a limit still has just one
nonconstant component, birational onto its image, by the same
degree argument.  It has no contracted components: any contracted
subtree would have a leaf of valency one and no markings, contrary
to stability.  Hence every unmarked limit has smooth domain
$\mathbb P^1$ and is birational onto its image.  Smooth-domain
birational stable maps represent the unparameterized birational-map
quotient underlying $\operatorname{RatCurves}(X)$.  The full
degree-$\iota$ stable-map locus is proper and has no boundary, so the
corresponding full irreducible component of this quotient is proper;
its finite normalization $V^j$ is proper as well.
These are actual proper irreducible components, as required by
Theorem~\ref{thm:ordered-chain-bound}.  Their numerical classes are
the independent classes $\gamma_j$, and the curves already selected
give the required chain.
\end{proof}

\begin{proposition}[Rigidity at equality]
\label{prop:equality-rigidity}
Suppose $\iota\ge2$, $r(\iota-1)=n$, and the independent nonzero
product in \eqref{eq:nonzero-independent-product} exists.  Then
\[
 X\cong(\mathbb P^{\iota-1})^r.
\]
\end{proposition}

\begin{proof}
Since each $d_{\gamma_j}\ge\iota$, equality in the grading bound
forces $d_{\gamma_j}=\iota$ for every $j$.
Lemma~\ref{lem:minimum-degree-incidence} supplies numerically
independent proper components $V^1,\ldots,V^r$ and curves
$C_1,\ldots,C_r$ forming a chain in that order.  Choose $x\in C_1$.
Theorem~\ref{thm:ordered-chain-bound} gives
\[
 \dim\operatorname{Locus}(V^1,\ldots,V^r)_x
 \ge r(\iota-1)=n.
\]
This locus is contained in the image of the universal evaluation of
$V^r$, so $V^r$ is covering.  Since $V^r$ is proper, its evaluation
image is closed; consequently every point of $X$ lies on a member
of $V^r$.

If $r>1$, choose a member of $V^r$ through a point of $C_1$ and
follow it by $C_1,\ldots,C_{r-1}$.  This is a chain in the order
$V^r,V^1,\ldots,V^{r-1}$.  Applying the same dimension bound proves
that $V^{r-1}$ is covering.  Repeating this cyclic rotation makes
each family the last member of a nonempty ordered chain, so every
$V^j$ is covering.  For $r=1$ the first application already suffices.

All hypotheses of Theorem~\ref{thm:occhetta-product} now hold with
$n_j=\iota-1$: the families are unsplit and covering, their
numerical classes are independent, their degrees are $\iota$, and
their contributions $n_j$ sum to $n$.  The resulting product is
$(\mathbb P^{\iota-1})^r$.
\end{proof}

Conversely, $(\mathbb P^{\iota-1})^r$ has dimension $r(\iota-1)$
and Picard number $r$.  Its anticanonical divisor has degree $\iota$
on a line in one factor, and degree at least $\iota$ on every
irreducible rational curve.  Its pseudoindex is therefore $\iota$,
and equality holds.
